% % % % % % % % % % % % % % % % % % % % % % % % % % % % % % % % % % % % % % % % % % % % % % % % % % % % % % % % % % % % % % % %
% SAT_STYLE_MATH_MCQS.tex
% 1_13.tex
% 
% encoding: UTF-8 
% EOL: LF
%
% format: LaTeX
% indent: spaces (2)
% width: 127
% % % % % % % % % % % % % % % % % % % % % % % % % % % % % % % % % % % % % % % % % % % % % % % % % % % % % % % % % % % % % % % %
Une montre à aiguilles indique $12$\,h\,$5$: combien vaut l'angle, arrondi sur deux décimales après la virgule, entre %
l'aiguille des heures et celle des minutes?
%
\begin{enumerate}
  \item $28,33°$
  \item $27,50°$
  \item $26,83°$
  \item $25,66°$
  \item Aucune des propositions ci-dessus
\end{enumerate}
% 
\begin{proof}
%
À $12h$ précises, cet angle vaut $0°$ car les aiguilles sont superposées. Cinq minutes plus tard, soit un douzième d'heure, %
l'aiguille des minutes a avancé d'exactement $\frac{360°}{12} = 30°$. Pendant ces mêmes cinq minutes, l'aiguille des heures a %
avancé dans le même sens, mais douze fois moins vite, de $30°/12 = 2,5°$ exactement. L'angle entre les deux aiguilles est %
donc de %
%
\begin{equation}
  30° - 2,5° = 27,5°.
\end{equation}
%
Réponse \textbf{B}.
\end{proof}